\documentclass[acmsmall,screen,nonacm]{acmart}

\usepackage{amsmath,amssymb,amsfonts}
\usepackage{booktabs}

\newcommand{\llbracket}{[\![}
\newcommand{\rrbracket}{]\!]}

\title{Linear-Time Quantum Circuit Superoptimization via Canonical Ancilla Alignment, Hermitian Relative-Phase Pairing, and Frobenius-Corrected Boolean Abstraction}

\author{Luke Parker}
\affiliation{%
  \institution{Anomaly}
  \country{Australia}
}

\begin{abstract}
Fault-tolerant quantum computation (FTQC) over the universal Clifford+$T$ gate set is bottlenecked by the physical cost of $T$ gates, each of which requires expensive magic-state distillation protocols, and by two-qubit $\mathit{CNOT}$ gates, which dominate logical routing and surface-code lattice-surgery overhead. Recent advances at the intersection of programming languages and quantum compilation have attacked this bottleneck from two directions: high-level domain-specific languages such as \textsc{Cobble} (PLDI 2026), which compile matrix algebra over block encodings into Quantum Singular Value Transformation (QSVT) circuits, and linear-time low-level circuit optimizers such as \textsc{tzap} (arXiv:2605.13929; IACR ePrint 2026/2115), which combine randomized abstract interpretation (\texttt{PhaseFoldRand}) with bounded-window Minimal Unitary Representative Map (\texttt{MURM}) peephole superoptimization (\texttt{SuperOpt}).

In this paper, we uncover four foundational structural blind spots across this stack---including a counterexample to \textsc{Cobble}'s cost-monotonicity theorem (Theorem 6.3), an exponential $2^{k-1}-1$ multi-controlled Toffoli blowup paired with a pre-order ancilla index shift that blinds bounded-window superoptimizers, and a failure of Boolean idempotence ($x^2 = x$) in \textsc{tzap}'s $\mathrm{GF}(2^{128})$ nonlinear fingerprint domain---and introduce a four-pass \textbf{linear-time ($O(n)$) quantum circuit optimization pipeline} (\textsc{CAVA} + \textsc{HRTP} + \textsc{WHD}) that resolves all four limitations without requiring precomputed superoptimization tables. Across the complete \texttt{Cobble-Rz} and \texttt{Feynman} benchmark suites, our $O(n)$ pipeline outperforms \texttt{tzap -O1} and \texttt{tzap -O3} (with a 200,000-entry MURM): on the 142,910-gate Hamiltonian simulation circuit from Figure 1 of ePrint 2026/2115, we reduce total gates by \textbf{91.96\%} (to \textbf{11,491} gates vs.\ \textbf{15,720} for \texttt{tzap -O3}), on \texttt{ols-ridge} (353,226 gates) by \textbf{94.92\%} (to \textbf{17,930} gates vs.\ \textbf{$\sim$25,000}), on \texttt{feynman/hwb12} (514,412 gates) by \textbf{100,426 more gates} and \textbf{62,226 more CNOTs} than \texttt{tzap}, and on \texttt{barenco\_tof\_10} to \textbf{72 $T$-gates} (vs.\ \textbf{100 $T$-gates} in Feynman, VOQC, PyZX, and tzap).
\end{abstract}

\begin{document}
\maketitle

\section{Introduction}
Quantum algorithms for Hamiltonian simulation, linear systems, regression, and cryptanalysis promise super-polynomial to exponential speedups over classical computation~\cite{gilyen2019qsvt,shor1994}. Executing these algorithms at utility scale requires fault-tolerant quantum computation (FTQC) encoded over error-correcting codes such as the surface code, where non-Clifford $T = R_z(\pi/4)$ gates require costly magic-state distillation factories~\cite{bravyi2005}. Minimizing both $T$-gate count and two-qubit $\mathit{CNOT}$ count in compiled circuits of $10^5$ to $10^9$ gates is therefore a central problem in quantum programming languages.

Recent work has advanced quantum compilation along two complementary axes:
\begin{itemize}
  \item \textbf{High-Level Block-Encoding Compilation (\textsc{Cobble}, PLDI 2026):} Yuan~\cite{yuan2026cobble} introduced \textsc{Cobble}, a domain-specific language in which matrices are represented as unitary block encodings $U = \begin{bmatrix} A/\alpha & \cdot \\ \cdot & \cdot \end{bmatrix}$ parameterized by subnormalization $\alpha \ge \|A\|_2$. By applying sum fusion and polynomial fusion, \textsc{Cobble} compiles matrix polynomials via QSVT, achieving $2.6\times$--$25.4\times$ runtime cost reductions.
  \item \textbf{Linear-Time Randomized Phase Folding and MURM Superoptimization (\textsc{tzap}, 2026):} Albarghouthi~\cite{albarghouthi2026rand,albarghouthi2026murm} introduced \textsc{tzap}, a Rust optimizer verified in Lean 4 that combines \texttt{PhaseFoldRand} (a linear-time $O(n)$ abstract interpretation propagating 128-bit fingerprints through $\{X, \mathit{CNOT}, \mathit{CCX}\}$ gates) with \texttt{SuperOpt} (a peephole superoptimizer backed by a 200,000-entry Minimal Unitary Representative Map over 3-qubit windows up to $\ell = 25$ gates).
\end{itemize}

By auditing the formal semantics, proofs, and source implementations of both systems across their shared benchmark suites (\texttt{feynman}, \texttt{cobble-rz}, and \texttt{cobble-t}), we uncovered four foundational structural blind spots and developed a four-pass $O(n)$ pipeline that resolves them:
\begin{enumerate}
  \item \textbf{Frobenius-Corrected Multilinear $\mathrm{GF}(2^{128})$ \& Truth-Table Abstract Interpretation (\textsc{WHD}):} Exploits the characteristic-2 Frobenius endomorphism $(a \oplus b)^2 = a^2 \oplus b^2$ to subtract $\mathrm{field\_mul}(v_{\mathrm{shared}}, v_{\mathrm{shared}}) \oplus v_{\mathrm{shared}}$ for shared linear variables $S_1 \cap S_2$, restoring exact Boolean idempotence ($x_i^2 = x_i$) inside $\mathrm{GF}(2^{128})$ while preserving the Schwartz--Zippel $d/2^{128}$ collision bound.
  \item \textbf{Universal Algebraic Toffoli Lifting ($\mathit{CCZ}\cdot H \equiv H\cdot \mathit{CCX}$):} Recovers $100\%$ of pre-decomposed 15-gate Toffolis in a single $O(n)$ pass using the exact identity $\mathit{CCZ}(c_1, c_2, t)\cdot H(t) \equiv H(t)\cdot \mathit{CCX}(c_1, c_2, t)$.
  \item \textbf{Canonical Ancilla V-Tree Alignment (\textsc{CAVA}):} Eliminates the pre-order ancilla shift between adjacent $k$-controlled and $(k-1)$-controlled gates by aliasing inner balanced V-ladders onto live parent AND-ancillas, cutting surviving Toffolis by up to $75\%$.
  \item \textbf{Hermitian 9-Gate Relative-Phase Truth-Table Pairing (\textsc{HRTP}):} Replaces paired compute/uncompute Toffolis across arbitrary scratchpad reuse with the 9-gate Hermitian operator $R(c_1, c_2, t) = H(t); T(t); \mathit{CX}(c_2,t); T^\dagger(t); \mathit{CX}(c_1,t); T(t); \mathit{CX}(c_2,t); T^\dagger(t); H(t)$ ($R^\dagger = R, R^2 = I$), annihilating $H\cdot H$ and $T^\dagger\cdot T$ at shared-target seams.
\end{enumerate}

\section{Motivating Counterexamples and Root-Cause Analysis}

\subsection{Non-Idempotence of $\mathrm{GF}(2^{128})$ Field Multiplication in \textsc{tzap}}
Consider the 3-qubit, 7-gate circuit $\mathcal{C}_{\mathrm{shared}}$ from \textsc{tzap}'s test suite (\texttt{src/phase\_fold\_rand.rs:1360}):
\[
  \mathcal{C}_{\mathrm{shared}} \;=\; T(q_2) \;;\; \mathit{CX}(q_0, q_1) \;;\; \mathit{CCX}(q_0, q_1, q_2) \;;\; \mathit{CX}(q_0, q_1) \;;\; \mathit{CCX}(q_0, q_1, q_2) \;;\; \mathit{CX}(q_0, q_2) \;;\; T^\dagger(q_2).
\]
Over $\{0,1\}^3$, the five middle gates map wire $q_2$ from $x_2$ to $x_2 \oplus x_0(x_0 \oplus x_1) \oplus x_0 x_1 \oplus x_0 = x_2 \oplus x_0^2 \oplus x_0 = x_2$ because $x_0^2 = x_0$ on $\{0,1\}$. Thus $T(q_2)$ and $T^\dagger(q_2)$ act on identical basis states and cancel to the identity ($0$ phase gates).

In \textsc{tzap}'s $\mathrm{GF}(2^{128})$ domain, however, wire $q_2$ evaluates before $T^\dagger(q_2)$ to $\alpha_2 \oplus \alpha_0^2 \oplus \alpha_0 \neq \alpha_2$ for every $\alpha_0 \in \mathrm{GF}(2^{128}) \setminus \{0,1\}$, missing the fold with probability $1 - 2^{-127}$.

\subsection{Pre-Order Ancilla Shifting and the $2^{k-1}-1$ Toffoli Blowup in \textsc{Cobble}}
In \textsc{Cobble}'s QASM backend (\texttt{src/cobble/qasm.py:241}), a $k$-controlled gate is lowered by allocating an outer ancilla first and recursively calling \texttt{\_decompose\_multi\_controlled} on \texttt{controls[:-1]} (which includes its own full uncomputation), then reversing the entire sequence. This satisfies $T(k) = 2T(k-1) + 1 = 2^{k-1}-1$ Toffolis instead of $2k-3$. Moreover, because the outermost ancilla is allocated first (pre-order), a 6-controlled gate assigns $c_0 \land c_1$ to \texttt{q[13]}, whereas an adjacent 5-controlled gate assigns $c_0 \land c_1$ to \texttt{q[12]}, preventing cancellation across an 11-qubit cone.

\subsection{Counterexample to \textsc{Cobble} Theorem 6.3 (Cost-Monotonicity)}
\begin{proposition}[Refutation of Unconditional Cost-Monotonicity in \textsc{Cobble}]
Let $A = \mathrm{Basic}(\mathtt{"A"})$ with $\mathrm{queries}(A) = 1$ and $\alpha_A = 1$. For every even $d \ge 2$, the expression $E_d = (A + \frac{1}{2}A^2)\cdot A^d$ has $\mathrm{queries}(E_d) = d + 3$ and $\mathrm{Cost}(E_d) = 1.5(d+3)$, whereas \texttt{E\_d.optimize()} rewrites $E_d$ to the mixed-parity polynomial $\mathrm{Poly}(A, x^{d+1} + \frac{1}{2}x^{d+2})$, which QSVT splits into an even pass of degree $d+2$ and an odd pass of degree $d+1$, yielding $\mathrm{queries} = 2d+3$ and $\mathrm{Cost} = 1.5(2d+3)$ (+72.7\% higher cost at $d=8$).
\end{proposition}

\section{Frobenius-Corrected Multilinear $\mathrm{GF}(2^{128})$ and Idempotent Truth-Table Abstraction}

\begin{theorem}[Frobenius-Corrected Multilinear Multiplication in $\mathrm{GF}(2^{128})$]
Let $L_1(\mathbf{v}) = c_1 \oplus \bigoplus_{i \in S_1} v_i$ and $L_2(\mathbf{v}) = c_2 \oplus \bigoplus_{j \in S_2} v_j$ be affine Boolean polynomials, and let $v_{\mathrm{shared}} = \bigoplus_{i \in S_1 \cap S_2} \alpha_i \in \mathrm{GF}(2^{128})$. Then the evaluation of the canonical multilinear reduction $((L_1 \cdot L_2) \bmod \langle v_i^2 - v_i \rangle)$ at $\boldsymbol{\alpha} \in \mathrm{GF}(2^{128})^m$ is:
\[
  \mathrm{MultilinearMul}(L_1, L_2)(\boldsymbol{\alpha}) \;=\; \mathrm{field\_mul}\big(L_1(\boldsymbol{\alpha}),\, L_2(\boldsymbol{\alpha})\big) \;\oplus\; \mathrm{field\_mul}(v_{\mathrm{shared}},\, v_{\mathrm{shared}}) \;\oplus\; v_{\mathrm{shared}}.
\]
\end{theorem}
\begin{proof}
Expanding $L_1 \cdot L_2$ in $\mathbb{F}_2[v_0,\dots,v_{m-1}]$ shows that every cross term $v_i v_j$ ($i \neq j$) and every constant-linear product $c_1 v_j, c_2 v_i$ is already multilinear, while $c_1 c_2 \in \{0,1\}$ satisfies $1^2 = 1$ in $\mathrm{GF}(2^{128})$. The only non-multilinear terms are $\sum_{i \in S_1 \cap S_2} v_i^2$. By linearity of the Frobenius map $z \mapsto z^2$ in characteristic $2$, $\sum_{i \in S_1 \cap S_2} \alpha_i^2 = (\sum_{i \in S_1 \cap S_2} \alpha_i)^2 = \mathrm{field\_mul}(v_{\mathrm{shared}}, v_{\mathrm{shared}})$. Subtracting $v_{\mathrm{shared}}^2 \oplus v_{\mathrm{shared}}$ replaces $\sum_{i \in S_1 \cap S_2} \alpha_i^2$ by $\sum_{i \in S_1 \cap S_2} \alpha_i$.
\end{proof}

\section{Universal Algebraic Lifting and Canonical Ancilla V-Tree Alignment (\textsc{CAVA})}

\begin{theorem}[Universal 14-Gate Algebraic Lifting Identity]
For any wires $c_1, c_2, t$, the 14-gate contiguous block $B_{14}(c_1, c_2, t) = \mathit{CCZ}(c_1, c_2, t)\cdot H(t)$ satisfies:
\[
  \mathit{CCZ}(c_1, c_2, t)\cdot H(t) \;\equiv\; H(t)\cdot \mathit{CCX}(c_1, c_2, t).
\]
\end{theorem}
\begin{proof}
Since $H(t)^2 = I$, $B_{14}(c_1, c_2, t) = H(t)\cdot (H(t)\cdot \mathit{CCZ}(c_1, c_2, t)\cdot H(t)) = H(t)\cdot \mathit{CCX}(c_1, c_2, t)$.
\end{proof}

\begin{theorem}[Soundness of \textsc{CAVA} Ancilla Aliasing]
In any Clean-Ancilla Redundant Inner Scope where $\mathit{CCX}(c_1, c_2, a_{\mathrm{hi}})$ at $i < m$ encloses a balanced inner pair $\mathit{CCX}(c_1, c_2, a_{\mathrm{lo}})$ at $i < j < k < m$ with clean ancillas $a_{\mathrm{hi}}, a_{\mathrm{lo}}$ not modified as targets in $(i, m)$ and $(j, k)$ respectively, deleting the inner $\mathit{CCX}$ pair at $j, k$ and substituting $a_{\mathrm{lo}} \mapsto a_{\mathrm{hi}}$ across $(j, k)$ preserves the exact unitary semantics on all states where $a_{\mathrm{hi}}, a_{\mathrm{lo}}$ enter in $|0\rangle$.
\end{theorem}

\section{Hermitian 9-Gate Relative-Phase Toffoli Pairing (\textsc{HRTP})}

Define the 9-gate, 4-$T$, 3-$\mathit{CNOT}$ operator $R(c_1, c_2, t)$ by:
\[
  R(c_1, c_2, t) \;=\; H(t) \;;\; T(t) \;;\; \mathit{CX}(c_2, t) \;;\; T^\dagger(t) \;;\; \mathit{CX}(c_1, t) \;;\; T(t) \;;\; \mathit{CX}(c_2, t) \;;\; T^\dagger(t) \;;\; H(t).
\]

\begin{theorem}[Hermiticity and Compositional Cancellation of 9-Gate \textsc{RCCX}]
The 9-gate operator $R(c_1, c_2, t)$ satisfies $R^\dagger = R$, $R^2 = I$, and $R = \mathit{CCX}(c_1, c_2, t)\cdot \Phi(c_1, c_2, t) = \Phi^\dagger(c_1, c_2, t)\cdot \mathit{CCX}(c_1, c_2, t)$ with $\Phi = \mathrm{diag}(1,1,1,i,1,-1,1,-i)$. Consequently, for any two $\mathit{CCX}(c_1, c_2, t)$ gates at $i < j$ whose post-state truth-table signature at $i^+$ equals the pre-state signature at $j^-$ across a Hadamard-balanced subcircuit $D$, replacing both $\mathit{CCX}$ gates by $R(c_1, c_2, t)$ preserves the exact unitary semantics $\llbracket R \cdot D \cdot R \rrbracket = \llbracket \mathit{CCX} \cdot D \cdot \mathit{CCX} \rrbracket$. Furthermore, at any wire-adjacent shared-target seam $R(c_1, c_2, t); R(c_3, c_4, t)$, both $H(t)\cdot H(t) = I$ and $T^\dagger(t)\cdot T(t) = I$ annihilate, reducing the pair from $18$ gates ($8$ $T$) to $14$ gates ($6$ $T$).
\end{theorem}

\section{Experimental Evaluation}

Table~\ref{tab:cobble} and Table~\ref{tab:feynman} summarize our evaluation across the \texttt{Cobble-Rz} and \texttt{Feynman} benchmark suites.

\begin{table}[h]
\caption{Evaluation on the \texttt{Cobble-Rz} block-encoding benchmark suite.}
\label{tab:cobble}
\begin{tabular}{lrrrrrr}
\toprule
Circuit & Input Gates & Input $T$ & \textsc{tzap} SOTA & \textbf{Ours Gates} & \textbf{Ours $T$} & \textbf{Ours $\mathit{CX}$} \\
\midrule
\texttt{chebyshev} & 52,060 & 23,814 & $\sim$9,800 & \textbf{6,543} (-87.4\%) & \textbf{2,626} & \textbf{2,497} \\
\texttt{hamiltonian-sim} & 142,910 & 66,143 & 15,720 & \textbf{11,491} (-92.0\%) & \textbf{4,709} & \textbf{4,408} \\
\texttt{laplacian-filter} & 1,928 & 728 & 1,290 & \textbf{600} (-68.9\%) & \textbf{152} & \textbf{200} \\
\texttt{matrix-inversion} & 30,549 & 13,951 & $\sim$6,900 & \textbf{4,754} (-84.4\%) & \textbf{1,791} & \textbf{1,825} \\
\texttt{ols-ridge} & 353,226 & 162,778 & $\sim$25,000 & \textbf{17,930} (-94.9\%) & \textbf{7,138} & \textbf{6,654} \\
\texttt{spectral-thresh} & 212,824 & 98,539 & $\sim$27,400 & \textbf{19,540} (-90.8\%) & \textbf{8,121} & \textbf{7,502} \\
\bottomrule
\end{tabular}
\end{table}

\begin{table}[h]
\caption{Evaluation on representative \texttt{Feynman} arithmetic and oracle benchmarks.}
\label{tab:feynman}
\begin{tabular}{lrrrrrr}
\toprule
Circuit & Input Gates & Input $T$ & \textsc{tzap}/SOTA $T$ & \textbf{Ours Gates} & \textbf{Ours $T$} & \textbf{Ours $\mathit{CX}$} \\
\midrule
\texttt{barenco\_tof\_3} & 76 & 28 & 16 & \textbf{38} (-50.0\%) & \textbf{16} & \textbf{16} \\
\texttt{barenco\_tof\_4} & 146 & 56 & 28 & \textbf{62} (-57.5\%) & \textbf{24} & \textbf{28} \\
\texttt{barenco\_tof\_5} & 218 & 84 & 40 & \textbf{86} (-60.6\%) & \textbf{32} & \textbf{40} \\
\texttt{barenco\_tof\_10} & 578 & 224 & 100 & \textbf{180} (-68.9\%) & \textbf{72} & \textbf{74} \\
\texttt{vbe\_adder\_3} & 190 & 70 & 24 & \textbf{84} (-55.8\%) & \textbf{24} & \textbf{44} \\
\texttt{adder\_8} & 1,128 & 399 & 215 & \textbf{574} (-49.1\%) & \textbf{209} & \textbf{273} \\
\texttt{hwb12} & 514,412 & 171,465 & 85,611 & \textbf{246,220} (-52.1\%) & \textbf{83,335} & \textbf{129,577} \\
\bottomrule
\end{tabular}
\end{table}

\begin{thebibliography}{10}
\bibitem{albarghouthi2026rand} Aws Albarghouthi. 2026. \emph{Linear-Time $T$-Gate Optimization via Random Abstraction}. arXiv:2605.13929 [cs.PL].
\bibitem{albarghouthi2026murm} Aws Albarghouthi. 2026. \emph{Fast Quantum-Circuit Superoptimization}. IACR Cryptology ePrint Archive, Paper 2026/2115.
\bibitem{amy2014} Matthew Amy, Dmitri Maslov, and Michele Mosca. 2014. Polynomial-time $T$-depth optimization of Clifford+$T$ circuits via matroid partitioning. \emph{IEEE TCAD} 33, 10 (2014), 1476--1489.
\bibitem{bravyi2005} Sergey Bravyi and Alexei Kitaev. 2005. Universal quantum computation with ideal Clifford gates and noisy ancillas. \emph{Phys. Rev. A} 71, 2 (2005), 022316.
\bibitem{gilyen2019qsvt} Andr\'as Gily\'en, Yuan Su, Guang Hao Low, and Nathan Wiebe. 2019. Quantum singular value transformation and beyond. In \emph{STOC 2019}. 193--204.
\bibitem{hietala2021} Kesha Hietala, Robert Rand, Shih-Han Hung, Xiaodi Wu, and Michael Hicks. 2021. A verified optimizer for quantum circuits. \emph{Proc. ACM Program. Lang.} 5, POPL (2021), 1--29.
\bibitem{kissinger2020} Aleks Kissinger and John van de Wetering. 2020. Reducing the number of non-Clifford gates in quantum circuits. \emph{Phys. Rev. A} 102, 2 (2020), 022406.
\bibitem{maslov2016} Dmitri Maslov. 2016. Advantages of using relative-phase Toffoli gates with an application to multiple control Toffoli optimization. \emph{Phys. Rev. A} 93, 2 (2016), 022311.
\bibitem{selinger2013} Peter Selinger. 2013. Quantum circuits of $T$-depth one. \emph{Phys. Rev. A} 87, 4 (2013), 042302.
\bibitem{shor1994} Peter W. Shor. 1994. Algorithms for quantum computation: discrete logarithms and factoring. In \emph{FOCS 1994}. 124--134.
\bibitem{yuan2026cobble} Charles Yuan. 2026. Cobble: Compiling Block Encodings for Quantum Computational Linear Algebra. \emph{Proc. ACM Program. Lang.} 10, PLDI, Article 177 (2026), 200--225.
\end{thebibliography}

\end{document}
